\section*{ADR-002: Dual Ensemble Architecture (DEA)}
\addcontentsline{toc}{section}{ADR-002: Dual Ensemble Architecture (DEA)}

\begin{description}[leftmargin=3cm,labelwidth=2.5cm]
\item[\textbf{Status}] Accepted
\item[\textbf{Date}] April 2024
\item[\textbf{Authors}] Symphony Architecture Team
\end{description}

\subsection*{Summary}

We adopted a Dual Ensemble Architecture (DEA) that separates extensions into two distinct layers: The Pit (Infrastructure as Extensions - IaE) for performance-critical operations running in-process, and The Grand Stage (User-Faced Extensions - UFE) for community-driven functionality running out-of-process. This architectural pattern enables both high performance and safe extensibility.

\subsection*{Context}

Symphony's extension system must balance competing requirements:

\begin{itemize}
\item Ultra-low latency for infrastructure operations (microsecond-level response times)
\item Safety and isolation for community-developed extensions
\item Scalability to support thousands of concurrent extensions
\item Reliability where extension failures don't crash the core system
\item Developer experience enabling rapid extension development and deployment
\end{itemize}

Traditional IDE architectures either sacrifice performance for safety (all extensions out-of-process) or safety for performance (all extensions in-process). Symphony's AI-first approach requires both characteristics simultaneously.

\subsection*{Decision}

We implemented a Dual Ensemble Architecture with two distinct extension execution environments:

\subsubsection*{The Pit (Infrastructure as Extensions - IaE)}
Five privileged Rust-powered extensions running in-process with the Conductor:
\begin{itemize}
\item Pool Manager: AI model lifecycle and resource management
\item DAG Tracker: Workflow dependency mapping and execution
\item Artifact Store: Institutional memory with versioning and quality scoring
\item Arbitration Engine: Intelligent resource allocation and conflict resolution
\item Stale Manager: Data lifecycle management and system optimization
\end{itemize}

\subsubsection*{The Grand Stage (User-Faced Extensions - UFE)}
Community-driven extensions running out-of-process with IPC communication:
\begin{itemize}
\item Instruments: AI/ML models as extensions
\item Operators: Workflow utilities and data processing
\item Motifs: UI enhancements and specialized editors
\end{itemize}

\subsection*{Rationale}

\subsubsection*{Performance Optimization}
The Pit's in-process execution enables microsecond-level response times for infrastructure operations. Benchmarking showed 100-1000× performance improvement over IPC-based alternatives for critical path operations.

\subsubsection*{Safety Through Isolation}
The Grand Stage's out-of-process execution provides complete isolation for community extensions. Extension crashes, memory leaks, or security vulnerabilities cannot affect the core system or other extensions.

\subsubsection*{Architectural Clarity}
Clear separation between infrastructure (The Pit) and user functionality (The Grand Stage) simplifies system reasoning, debugging, and maintenance. Each layer has distinct performance characteristics and development constraints.

\subsubsection*{Scalability Strategy}
The architecture scales differently for each layer: The Pit maintains fixed, optimized components while The Grand Stage can scale horizontally with additional processes and resources.

\subsubsection*{Development Experience}
Developers can choose the appropriate execution model based on their requirements: high-performance infrastructure extensions (The Pit) or safe, flexible user extensions (The Grand Stage).

\subsubsection*{Trade-offs Accepted}
\begin{itemize}
\item Increased architectural complexity compared to single-layer approaches
\item Different development models for each layer requiring distinct expertise
\item IPC overhead for Grand Stage extensions (0.1-0.5ms latency)
\item Limited number of Pit extensions to maintain performance guarantees
\end{itemize}

\subsection*{Alternatives Considered}

\subsubsection*{Monolithic In-Process Architecture}
\textbf{Pros}: Maximum performance, simple communication, unified development model
\textbf{Cons}: Extension crashes affect entire system, security vulnerabilities, difficult debugging, limited scalability
\textbf{Rejected because}: Safety and reliability requirements outweigh performance benefits for user-facing functionality

\subsubsection*{Pure Out-of-Process Architecture}
\textbf{Pros}: Complete isolation, excellent security, independent scaling, simplified debugging
\textbf{Cons}: IPC overhead for all operations, complex state management, performance bottlenecks
\textbf{Rejected because}: Performance requirements for infrastructure operations cannot be met with IPC overhead

\subsubsection*{Microservice Architecture}
\textbf{Pros}: Independent deployment, technology diversity, horizontal scaling, fault isolation
\textbf{Cons}: Network latency, complex service discovery, distributed system complexity, resource overhead
\textbf{Rejected because}: Network latency incompatible with microsecond-level performance requirements

\subsection*{Consequences}

\subsubsection*{Positive}
\begin{itemize}
\item Achieved target performance: 10-50 microseconds for Pit operations, 0.1-0.5ms for Grand Stage
\item Zero system crashes from extension failures over 6 months of operation
\item Successful community adoption with 200+ extensions developed for The Grand Stage
\item Clear architectural boundaries simplify system maintenance and debugging
\item Flexible scaling: Pit maintains consistent performance while Grand Stage scales with demand
\item Enhanced security through process isolation for untrusted community code
\end{itemize}

\subsubsection*{Negative}
\begin{itemize}
\item Increased development complexity requiring expertise in both execution models
\item IPC overhead limits some use cases for Grand Stage extensions
\item Resource overhead from multiple processes (50-100MB per Grand Stage extension)
\item Complex debugging scenarios spanning both execution environments
\item Limited flexibility for promoting Grand Stage extensions to Pit status
\end{itemize}

\subsection*{Success Criteria}

\begin{itemize}
\item Pit operations complete within 100 microseconds (Achieved: 10-50 microseconds)
\item Grand Stage IPC latency under 1ms (Achieved: 0.1-0.5ms)
\item Zero system crashes from extension failures (Achieved: 0 crashes)
\item Support 100+ concurrent Grand Stage extensions (Achieved: 200+ extensions)
\item Community adoption of extension development (Achieved: Active developer ecosystem)
\end{itemize}

\subsection*{Risks and Mitigations}

\begin{center}
\begin{tabular}{@{}p{4cm}p{2cm}p{6cm}@{}}
\toprule
\textbf{Risk} & \textbf{Impact} & \textbf{Mitigation} \\
\midrule
Architectural complexity & High & Comprehensive documentation and developer tooling \\
IPC performance bottlenecks & Medium & Optimized IPC implementation and caching strategies \\
Resource consumption & Medium & Process pooling and resource monitoring \\
Development model confusion & Medium & Clear guidelines and examples for each layer \\
Pit extension limitations & Low & Careful selection criteria and upgrade paths \\
\bottomrule
\end{tabular}
\end{center}

The Dual Ensemble Architecture successfully balances Symphony's competing requirements for performance and safety, enabling both microsecond-level infrastructure operations and safe community extensibility. This architectural pattern has become a key differentiator for Symphony's AI-first development environment.